\begin{center}
    {\LARGE \textbf{2010无机化学}} \\[2ex]
    {\large 时量180分钟 \quad 满分100分}
\end{center}

\vspace{0.5cm}
\noindent\textbf{注意事项:}
\begin{enumerate}[label=\arabic*., parsep=0pt, itemsep=0pt]
    \item 答题前,考生先将自己的姓名、学号、考场号、座位号填写在答题卡上,将条形码准确粘贴在考生信息条形码粘贴区。
    \item 考试结束后,将本试卷与答题纸一并收回。
    \item 回答各个问题时,将答案写在答题纸上,写在本试卷上无效。
\end{enumerate}

\vspace{0.5cm}
\noindent\textbf{一、完成并配平下列反应的方程式（24 pts.）}

\begin{enumerate}[label=\arabic*., start=1, itemsep=1.5ex]
    \item 在酸性\ce{(NH4)2S2O8}溶液中加入\ce{MnSO4}和几滴\ce{AgNO3}溶液。
    \item 将\ce{H2O}滴在红磷与\ce{I2}的混合物上。
    \item 将\ce{Cr2O3}与\ce{KHSO4}共熔。
    \item 用\ce{Zn}粉还原亚硫酸氢钠。
    \item 将\ce{Co2O3}溶于盐酸。
    \item 金溶于王水中。
\end{enumerate}

\vspace{0.5cm}
\noindent\textbf{二、合成与制备（18 pts.）}

\begin{enumerate}[label=\arabic*., start=1, itemsep=1.5ex]
    \item 以铬绿（\ce{Cr2O3}）为原料制备重铬酸钾，写出相关的反应方程式。
    \item 用方程式表示由\ce{Na2S}，\ce{Na2CO3}及\ce{SO2}为反应物合成硫代硫酸钠。
    \item 以海水为主要原料，使用\ce{Cl2}，\ce{Na2CO3}等试剂制备单质溴，写出相关的反应方程式。
\end{enumerate}

\vspace{0.5cm}
\noindent\textbf{三、鉴别与分离（22 pts.）}

\begin{enumerate}[label=\arabic*., start=1, itemsep=1.5ex]
    \item 现有三瓶无色溶液，可能为\ce{AsCl3}、\ce{SbCl3}和\ce{BiCl3}，试加以鉴别。
    \item 试用两种方法鉴别下列各物质（1）\ce{NaNO3}和\ce{NH4NO3}；（2）\ce{Fe^{2+}}和\ce{Fe^{3+}}
    \item 设计方案将含有\ce{Ag^{+}}、\ce{Zn^{2+}}、\ce{Cd^{2+}}、\ce{Hg^{2+}_{2}}的混合溶液分离。
\end{enumerate}

\vspace{0.5cm}
\noindent\textbf{四、回答下列各题（60 pts.）}

\begin{enumerate}[label=\arabic*., start=1, itemsep=1.5ex]
    \item 试说明$H$原子中$3s$和$3p$轨道能量相等，而$Cl$原子中$3s$轨道能量比相应的$3p$轨道能量低。
    \item 利用价层电子对互斥理论写出下列分子或离子的几何构型，并指出中心原子的杂化方式。
    
    {\centering
    \ce{POCl3}，\ce{SOCl2}，\ce{ICl2^{+}}，\ce{S2O3^{2-}}，\ce{SF4} \par
    }
    \item 比较下列各对物质熔沸点高低，并简述理由。
    \begin{enumerate}[label=(\arabic*), itemsep=1ex]
        \item \ce{NH3} 和 \ce{PH3}
        \item \ce{FeCl3} 和 \ce{FeCl2}
        \item \ce{Br2} 和 \ce{ICl}
    \end{enumerate}
    \item 预测下列各组配合物稳定性大小，并说明原因。
    \begin{enumerate}[label=(\arabic*), itemsep=1ex]
        \item \ce{[Co(NH3)6]^{3+}} 和 \ce{[Co(NH3)6]^{2+}}
        \item \ce{HgF4^{2-}}，\ce{HgCl4^{2-}}，\ce{HgBr4^{2-}}，\ce{HgI4^{2-}}
    \end{enumerate}
    \item 试解释下列现象：
    \begin{enumerate}[label=(\arabic*), itemsep=1ex]
        \item 浓\ce{H2SO4}和浓\ce{HClO4}具有较强的氧化性，而稀\ce{H2SO4}和稀\ce{HClO4}不具有氧化性。
        \item 向\ce{MnSO4}溶液中滴加\ce{NaOH}溶液生成白色沉淀，沉淀暴露在空气中逐渐变为棕黑色，将棕黑色沉淀用酸性\ce{H2O2}溶液处理，沉淀消失，得到无色溶液。
    \end{enumerate}
    \item 试说明为什么氟的电子亲和能比氯低，但\ce{F2}却比\ce{Cl2}活泼。
    \item 根据下列数据：
    \[
    \begin{array}{ll}
    E^\ominus(\ce{Ag+}/\ce{Cu+}) = 0.799\,V, & E^\ominus(\ce{Cu^{2+}}/\ce{Cu+}) = 0.153\,V, \\
    K_{SP}(\ce{AgCl}) = 1.8 \times 10^{-10}, & K_{SP}(\ce{CuCl}) = 1.2 \times 10^{-6}
    \end{array}
    \]
    计算电池反应 \ce{Ag + Cu^{2+} + 2Cl- \rightleftharpoons CuCl + AgCl} 的电动势 $E^\ominus$。
    \item 根据晶体场理论说明下列配合物的$d$电子排布并计算晶体场稳定化能。
    已知：
    \[
    \begin{array}{lcc}
    \hline
     & \ce{[Co(NH3)6]^{3+}} & \ce{[Fe(H2O)6]^{3+}} \\
    \hline
    \Delta/\text{cm}^{-1} & 22000 & 23000 \\
    P/\text{cm}^{-1} & 30000 & 13700 \\
    \hline
    \end{array}
    \]

\end{enumerate}

\vspace{0.5cm}
\noindent\textbf{四、推理判断（26 pts.）}

\begin{enumerate}[label=\arabic*., start=1, itemsep=1.5ex]
    \item 在一种含有配离子$A$的溶液中，加入稀硝酸，有刺激性气体$B$、黄色沉淀$C$和黑色沉淀$H$产生。气体$B$能使\ce{KMnO4}溶液褪色。若通氯气于溶液$A$中，得到白色沉淀$I$和含有$D$的溶液。$D$与\ce{BaCl2}作用，有不溶于酸的白色沉淀$E$产生。若在溶液$A$中加入$KI$溶液，产生黄色沉淀$F$，再加入\ce{NaCN}溶液，黄色沉淀$F$溶解，形成无色溶液$G$。向$G$中通入\ce{H2S}气体，得到黑色沉淀$H$。根据上述实验结果，试确定$A$、$B$、$C$、$D$、$E$、$F$、$G$、$H$及$I$各为何物。并写出由$A$生成$B$、$C$、$H$及由$A$生成$D$和$I$的化学反应方程式。
    \item 白色晶体$A$与氢氧化钠固体共热得无色气体$B$和白色固体$C$。将气体$B$与黑色氧化物$D$共热得到红色金属单质$E$和化学性质不活泼的无色气体$F$，将气体$B$通入氯化汞溶液有白色沉淀$G$生成。将$C$加入酸化后的碘化钾溶液中，则溶液变黄，说明有$H$生成。将$C$加入酸性高锰酸钾溶液中，则溶液的紫色褪去，$C$被氧化成化合物$I$。加热分解白色晶体$A$生成无色气体$F$和水。试给出$A$、$B$、$C$、$D$、$E$、$F$、$G$、$H$和$I$的化学式，并写出由$B$生成$G$及由$C$生成$I$的化学反应方程式。
\end{enumerate}